\section{Conclusions}\label{sec:conclusions}

On this agentic AutoML workload, the Qwen3 sparse proposer has a clear energy
advantage over the dense proposer. It uses 42.0\,\% less GPU-board energy per
session (90.7 vs.\ 156.4\,kJ). Reconstructing the archived AMD CPU-package
counter over the same session boundaries gives 179.8 vs.\ 294.5\,kJ, a 39.0\,\%
reduction at the broader measured-component scope. Neither quantity is idle-
subtracted, and the latter is not full-system energy because DRAM and other host
components are absent.

The MoE's mean test accuracy is 2.85\,pp higher (0.8224 vs.\ 0.7939), but its
95\,\% interval includes zero. We therefore conclude that the energy advantage is
established for this design while accuracy superiority and equivalence are not.
The MoE retained 17 improvements against 12 and had one no-progress session
against five. Correctly pooling all sessions gives 64.0 vs.\ 156.4\,kJ of GPU
energy per retained improvement ($-59.1\,\%$); adding CPU-package energy gives
126.9 vs.\ 294.5\,kJ ($-56.9\,\%$). No-progress sessions finish on the starting
recipe under keep/revert semantics, not below it.

Loop budget is the clearest agent-policy result. Raising the budget from 10 to 20
increases pooled GPU energy per retained improvement by 146.2\,\%, with a
bootstrap interval excluding no change. The corrected patience estimate changes
direction relative to the original per-session-ratio analysis and remains
uncertain. Of Phase-1 GPU energy, 78.1\,\% is attributable to non-retained
iterations, 8.2\,\% to retained iterations, and 13.6\,\% lies outside aligned
iteration windows. This is energy productivity under a fixed budget: the design
does not show that poor proposals cause the loop to schedule more training.

The sample Pareto frontier contains two MoE cells and one dense cell. The cheapest
MoE cell is the natural low-energy choice, but exact frontier membership on the
accuracy axis is provisional because cell means are noisy. The measured keep
threshold is a within-session validation rule, not a between-cell test-accuracy
equivalence margin.

The exploratory reasoning experiment suggests an architecture-by-reasoning
interaction: reasoning raises energy in both arms, improves dense accuracy in
this sample, and does not improve the MoE sample. MoE-without-reasoning reaches
nearly the same observed mean accuracy as dense-with-reasoning for roughly one
third of its GPU energy, but the accuracy interval is too wide for equivalence.
The temperature follow-ups establish that $T=0$ increases exact duplication on
both arms. They do not establish a null effect on retention or accuracy, because
each level has three sessions and the original restart changed prompt content and
temperature together.

Three extensions are now cleanly separated from completed evidence. First, a
100-iteration MoE session with post-hoc replay of every retained checkpoint is
implemented but not yet run; the manuscript will report it only after its raw
NVML and EnergiBridge traces, trajectory, and replay results are recovered.
Second, CPU-served proposer inference (with training still on GPU~0) awaits
supervisor feedback and is not part of the current conclusions. Third, wider
workload and model-family axes, including the original \texttt{nanochat} workload,
are needed before treating the observed frontier as workload-invariant.
